\scanpage{20}
$G$ groupe

$k_0(G)$ anneau des $G$-ens. finis virtuels

\deleted{\unsure{Suranneau}}
\[
k_n(G)\deleted{\unsure{(K_0)}}=k(\overbrace{G\times\mathfrak S_n}^{G_{(n)}})
\qquad(n\geq0;\ \text{on pose }\mathfrak S_0=\{1\})
\]
on le considère comme $k_0(G)$-module
\[
k_*(G)\deleted{\unsure{(K_0)}}=\deleted{\unsure{\prod}}\bigoplus_{n\geq0}k_n(G)\deleted{\unsure{(K_0)}}
\]
\margin{(somme de \deleted{groupes abéliens}\inserted{$k_0(G)$-modules}) (gradué)}

On définit dessus une structure de $k_0(G)$-\uline{algèbre},
\deleted{en}\inserted{en} définissant la multiplication $*$ par
\[
\operatorname{cl}_m(E)*\operatorname{cl}_n(F)=\operatorname{cl}_{m+n}((E\times F)\mathbin{\wedge}_{\mathfrak S_m\times\mathfrak S_n}\mathfrak S_{m+n})
\]
où $E$ est un $G_{(m)}$ ens. fini, $F$ un $G_{(n)}$ ens fini;
$E\times F$ est un $G\times G\times\mathfrak S_m\times\mathfrak S_n$ ens. fini,
$\mathfrak S_m\times\mathfrak S_n\longrightarrow\mathfrak S_{m+n}$ l'hom bien connu,
donc $(E\times F)\mathbin{\wedge}_{\mathfrak S_m\times\mathfrak S_n}\mathfrak S_{m+n}$ est un $G\times G\times\mathfrak S_{m+n}$
ensemble, qu'on considère comme un $G\times\mathfrak S_{m+n}$
ensemble grâce à l'application diagonale, ce qui
donne un sens à la formule.

Si $E$ est un $G$-ens. \unsure{finis|fini}, considérons
\[
\begin{array}{l}\deleted{\tau_t(E)=\sum\unsure{nt}}\\
\tau^n(E)=\operatorname{cl}_n(\unsure{\operatorname{Mon}|\operatorname{Hom}}(I_n,E))\in k_n(G)\deleted{\unsure{(K_0)}}\qquad(n\geq0)\end{array}
\]
où $I_n=[1,n]$, et
\[
\tau^*_{\deleted{\unsure{t}}}(E)\unsure{=}\sum_{n\geq0}\tau^n(E)\deleted{\unsure{\in}}\quad\deleted{\unsure{\in}}\in\uline{1+\widehat{k_*(G)}^+}
\]
(\uline{NB} $\tau^0(E)=1$, $\tau^1(E)=\deleted{E}\operatorname{cl}_0(E)$)

\uline{Prop} On a
\[
\tau^*(E\amalg F)=\tau^*(E).\tau^*(F)
\]
i.e.
\[
\tau^n(E\amalg F)=\sum_{i+j=n}\tau^i(E)\tau^j(F)
\]
\mspage{19}
